\section{Algebraic Identities}

\subsection{Introduction}

Algebraic identities are formulas that allow us to rewrite expressions in a
more useful or factored form. They play a central role in number theory,
especially when studying divisibility, prime numbers, and the structure of
integers.

In this chapter we explore classical identities, their applications, and
historical curiosities.

\subsection{Difference of Powers}

One of the most important identities is the factorization of a difference of
powers:

\begin{theorem}[Difference of Powers]
For any integer \(x\) and positive integer \(n\),


\[
x^n - 1 = (x - 1)(x^{n-1} + x^{n-2} + \cdots + x + 1).
\]


\end{theorem}

This identity is fundamental in the study of numbers of the form
\(x^n - 1\), especially when \(x = 2\).

\begin{example}[Factoring a Mersenne-Type Number]
Consider the number


\[
2^{15} - 1.
\]


Since \(15 = 3 \cdot 5\), we can apply the identity twice:


\[
2^{15} - 1 = (2^3 - 1)(2^{12} + 2^9 + 2^6 + 2^3 + 1).
\]


Then factor \(2^3 - 1\) again:


\[
2^3 - 1 = (2 - 1)(2^2 + 2 + 1) = 7.
\]


This method reveals the internal structure of the number without guessing
prime factors.
\end{example}

\subsection{Sum of Powers}

Another classical identity is the sum of powers:

\begin{theorem}[Sum of Powers]
For any integer \(x\) and positive integer \(n\),


\[
x^n + 1 =
\begin{cases}
(x + 1)(x^{n-1} - x^{n-2} + \cdots - x + 1), & \text{if } n \text{ is odd} \\
\\[6pt]
\text{irreducible over the integers}, & \text{if } n \text{ is even}.
\end{cases}
\]


\end{theorem}

\begin{example}[A Curious Identity]
Since \(5\) is odd,


\[
2^5 + 1 = 33 = (2 + 1)(2^4 - 2^3 + 2^2 - 2 + 1).
\]


This identity explains why some numbers of the form \(2^n + 1\) factor,
while others (the Fermat numbers) are suspected to be prime.
\end{example}

\subsection{Historical Notes}

\begin{remark}[Euclid and Mersenne]
Euclid already used the identity


\[
x^n - 1 = (x - 1)(x^{n-1} + \cdots + 1)
\]


in his proof that there are infinitely many primes.  
Much later, Marin Mersenne studied numbers of the form \(2^p - 1\), now called
\emph{Mersenne numbers}. Some of these are prime, and they produce the largest
known primes in history.
\end{remark}

\begin{remark}[Fermat's Curiosity]
Pierre de Fermat conjectured that all numbers of the form


\[
F_n = 2^{2^n} + 1
\]


were prime.  
He was correct for \(n = 0,1,2,3,4\), but Euler disproved the conjecture by
showing:


\[
F_5 = 2^{32} + 1 = 4294967297 = 641 \cdot 6700417.
\]


This is one of the most famous false conjectures in number theory.
\end{remark}

\subsection{Applications}

\begin{itemize}
    \item Factoring large integers using algebraic structure.
    \item Understanding prime candidates such as Mersenne and Fermat numbers.
    \item Simplifying expressions in modular arithmetic.
    \item Building conjectures about divisibility and primality.
\end{itemize}

\subsection{Curiosities}

\begin{itemize}
    \item The largest known prime numbers are almost always of the form
    \(2^p - 1\), where \(p\) is prime.
    \item The identity for \(x^n - 1\) is used in cryptography, especially in
    algorithms involving cyclic groups.
    \item Some identities, like the sum of powers, behave differently depending
    on whether the exponent is odd or even — a subtle but important detail.
\end{itemize}
